%%%%

%%%   Самарский государственный технический университет

%%%   Кафедра прикладной математики и информатики

%%%

%%%   Преамбула для статьи в журнал "Вестник Самарского государственного

%%%   технического университета. Серия: «Физико-математические науки»"

%%%   Ориентирована под MikTEx 2.4 for Win32

%%%

%%%   Хотя сам журнал собирается под GNU/Linux на дистрибутиве TeXLive



\documentclass[11pt,a4paper,twoside]{article}



\usepackage[cp1251]{inputenc}

\usepackage[T2A]{fontenc}

\usepackage[english,russian]{babel}

\usepackage{samgtubib}

\usepackage[dvips]{graphicx}

\usepackage[multiple]{footmisc}



\usepackage{samgtu}

\renewcommand{\VestnikYear}{2009} % Год издания Вестника

\renewcommand{\VestnikNumber}{2\,(19)} % Номер Вестника

\renewcommand{\VestnikMonth}{Ноябрь} % Месяц выхода Вестника

\renewcommand{\VestnikData}{01.11.09} % Дата подписания Вестника в печать

\renewcommand{\VestnikUPL}{26,64} % Усл. печ. л.

\renewcommand{\VestnikUIL}{25,73} % Уч.-изд. л.

\renewcommand{\VestnikRegNumber}{479/09} % Регистрационный номер

\renewcommand{\VestnikOrder}{994} % Номер заказа






%
% \begin{document}


%\renewcommand{\baselinestretch}{0.85}

%\ShortPartition{Математическое моделирование}{}

\renewcommand{\baselinestretch}{0.85}

\renewcommand{\firstpage}{199}
\renewcommand{\lastpage}{201}

\udk{517.958:531.72}
\msc{26A33; 28A80}% Коды MSC см. http://www.ams.org/msc/

\title{О математическом моделировании процесса фрактальной
миграции загрязнений в~природных пористых системах}% Полный заголовок
{О математическом моделировании процесса фрактальной миграции загрязнений\dots}% Заголовок, который идёт в верхний колонтитул

\author{А.\,А.~Вендина}{В\,е\,н\,д\,и\,н\,а~А.\,А.}

\institute{Северо-Кавказский государственный технический университет,\\
355029, Ставрополь, просп. Кулакова, 2.}


\email{E-mail}{Roven-ka@yandex.ru} %E-mail's авторов

\otherinf{\fullauthor{Алла Анатольевна Вендина} \position{ст. преподаватель} \dept{каф. высшей математики}}

\keywords{миграция загрязнённых подземных вод, диффузия во фрактальных средах, дробные дифференциальные уравнения}

\workabstract{Рассматриваются вопросы математического моделирования нелинейных 
процессов миграции загрязнённых подземных вод, которые не могут быть описаны в рамках классического 
подхода в теории массопереноса. Для исследования пространственно"=временных 
закономерностей нелинейных эффектов, обусловленных масштабной 
инвариантностью, корректно поставлена нелокальная краевая задача для модельного дифференциального
уравнения нелинейной миграции.}


\engtitle{On the mathematical modeling of the contaminated groundwater fractal migration in natural porous systems}

\engauthor{A.\,A.~Vendina}{V\,e\,n\,d\,i\,n\,a~A.\,A.}

\enginstitute{North Caucasus State Technical University, \\ 
2, prosp. Kulakova,  Stavropol', 355029, Russia.}

\otherenginf{\fullauthor{Alla A. Vendina} \position{Senior Teacher} \dept{Dept. of Higher Mathematics}}

\engabstract{The problems of the mathematical modeling of the contaminated groundwater migration, which cannot be described in the context of the mass transfer theory classical approach, are considered. For the study of space-time conformity of nonlinear effects, determined by scaled invariance, the well-set nonlocal boundary problem for the model differential equation of nonlinear migration is proposed.}

\engkeywords{migration of contaminated groundwater, diffusion on fractals, fractional differential equations}


\date{03/III/2011} % Дата поступления статьи в редакцию.
\revisiondate{07/VI/2011} % В окончательном виде

%%%%%%%%%%%%%%%%%%%%%%%%%%%%%%%%%%%%%%%%%%%%%%%%%%%%%%%%%%%%%%%%%%%%%%%%%%%%%%%%%%%%


\maketitle

\small


Процесс переноса загрязнений потоком подземных вод в~природных пористых средах в~рамках простых идеализированных 
представлений порового пространства и~справедливости принципа локальности при определённой схематизации удовлетворительно описывается системой дифференциальных уравнений \cite{ven:bib:1}, состоящей из 
\begin{itemize}
\item уравнения движения: 
\begin{equation}
\label{vend:eq1}
v=-k{\rm grad}\,I,
\quad
u_x=-v_x c+D^*\frac{\partial c}{\partial x},
\quad
u_y=-v_y c+D^*\frac{\partial c}{\partial y};
\end{equation}
\item уравнения сохранения:
\begin{equation}
\label{vend:eq2}
\frac{\partial(pv_x)}{\partial x}+\frac{\partial(pv_y)}{\partial y}=\frac{\partial(mp)}{\partial t},
\quad
m\frac{\partial c}{\partial t}=\frac{\partial u_x}{\partial x}+\frac{\partial u_y}{\partial y}-\frac{\partial N}{\partial t};
\end{equation}
\item уравнения состояния:
\begin{equation}
\label{vend:eq3}
q=q(p;p_c;\mu),
\quad
f=f(\rho;p),
\quad
\frac{\partial N}{\partial t}=\psi(N;c;\gamma_i).
\end{equation}
\end{itemize}
Здесь $v=v(x;y;t)$ "--- скорость движения потока грунтовых вод; $I=I(x;y;t)$ "--- напор; 
$c=c(x;y;t)$ "--- концентрация загрязняющих веществ в потоке грунтовых вод в~точке $(x;y)$ в~момент времени $t$; 
$m$ "--- активная пористость грунта; 
\mbox{$D^*=D^*(x;y;t)$} "--- коэффициент конвективной диффузии; 
$p=p(x;y;t)$ "--- давление; $p_c$ "--- давление пласта; $\mu$ "--- коэффициент упругости; $\rho$ "--- плотность жидкости; 
$N= N(x;y;t)$ "--- концентрация вещества в~твёрдой фазе; 
$\gamma _i$ "--- коэффициент скорости массообмена.

В силу уравнений \eqref{vend:eq1}--\eqref{vend:eq3}, учитывающих совместное действие гравитационных и диффузионных сил,  одномерный процесс загрязнения подземных вод
описывается нелинейным дифференциальным уравнением в частных производных 
второго порядка параболического типа:
\begin{equation}
\label{vend:eq4}
m\frac{\partial c}{\partial t}=\frac{\partial}{\partial x}\Bigl[D(x;t)\frac{\partial c}{\partial x}\Bigr]-\frac{\partial[v(x;t)c]}{\partial x}-\frac{\partial N}{\partial t}.
\end{equation}
Обобщённое дифференциальное уравнение \eqref{vend:eq4} в~рамках   эмпирических законов Фурье и~Фика описывает процесс изменения во времени  концентрации загрязнений в~какой"=либо точке среды вследствие разности 
переноса загрязнений фильтрационным потоком и гидродинамической дисперсии 
при наличии растворения и поступления в раствор солей и примесей, 
содержащихся в почвогрунтах. Построенные на основе обобщенного уравнения 
\eqref{vend:eq4} решения различных начально"=краевых задач [\citen{ven:bib:1}, \citen{ven:bib:2}] не учитывают влияния 
неоднородной структуры порового пространства на процессы миграции и, как 
следствие, не отражают возникающих в них нелинейных эффектов, например, таких 
как <<эффекты памяти>> и~конечность скорости распространения.

Современный метод математического моделирования возникающих в процессе миграции 
нелинейных явлений существенно 
использует идеи фрактальной геометрии в сочетании с формализмом дробного 
интегро"=диф\-фе\-рен\-ци\-ро\-ва\-ния. Почва и~почвогрунт хорошо интерпретируются как среды, структура пор 
которых образует множество с фрактальной размерностью $d_f$. 

Пусть за промежуток времени $\Delta t$ происходит изменение концентрации 
загрязнений подземных вод на величину 
$\delta c = {c}'_t \Delta t$.
В соответствии с методами фрактального анализа \cite{ven:bib:3} для скорости ${c}'_t $ изменения 
 концентрации загрязнения в~процессе нелинейной миграции 
принимается гипотеза о том, что она определяется как дробная 
производная в смысле Римана"--~Лиувилля концентрации $c(x;t)$:
\begin{equation}
\label{vend:eq5}
{c}'_t\approx D_{0t}^\alpha c(x;t)=\frac{1}{\Gamma(1-\alpha)}\frac{\partial}{\partial t}\int_0^t{\frac{c(x;\tau)d\tau}{(t-\tau)^\alpha}},
\end{equation}
где $\alpha\in(0; 1)$ "--- порядок дробной производной. Для 
широких классов пористых сред полагают $\alpha\approx d_f$.

При выполнении гипотезы \eqref{vend:eq5} обобщением уравнения \eqref{vend:eq4} является 
нелокальное дифференциальное уравнение в~частных производных второго порядка 
параболического типа с дробной производной по времени:
\begin{equation}
\label{vend:eq6}
mD_{0t}^\alpha c=\frac{\partial}{\partial x}\Bigl[D_\alpha(x;t)\frac{\partial c}{\partial x}\Bigr]-\frac{\partial[v_\alpha(x;t)c]}{\partial x}-\frac{\partial N}{\partial t},
\end{equation}
где $D_\alpha$ "--- коэффициент фрактальной диффузии; $v_\alpha $ "--- 
фрактальная скорость. 

Для однозначного описания нелинейной фрактальной миграции загрязнений к 
модельному уравнению миграции \eqref{vend:eq6} присоединяют \cite{ven:bib:4} видоизмененные начальные 
условия:  условие типа Коши
\begin{equation}
\label{vend:eq7}
\lim\limits_{t\to0}D_{0t}^{\alpha-1}c(x;t)=f(x)
\end{equation}
или локальное весовое начальное условие
$$\lim\limits_{t\to0}t^{1-\alpha}c(x;t)=f(x).$$

В условиях \eqref{vend:eq7} дробный показатель $\alpha$, отражая фрактальную размерность 
пористой среды, соответствует доле пористых каналов или ветвей, открытых для 
протекания загрязнённых подземных вод в начальный момент времени. При 
решении задач с прикладной направленностью использование начальных условий 
типа Коши \eqref{vend:eq7} оказывается непригодным для практических расчётов.  
Для преодоления указанной проблемы, как правило, дробная производная Римана"--~Лиувилля заменяется на нормализованную производную по Капуто и дифференциальное уравнение
рассматривается с классическим начальным условием Коши 
\begin{equation}
\label{vend:eq8}
c(x;0)=f(x).
\end{equation}

Для исследования асимптотического режима миграции загрязнённых 
подземных вод в пористых средах в настоящей работе рассматривается следующая нелокальная 
начально"=краевая задача.

\begin{newthm}{Задача}Найти в прямоугольной области $\Omega=\{(x;t): 0<x<l, 0<t<T\}$ регулярное решение $c=c(x;t)$ уравнения \eqref{vend:eq6}$,$ удовлетворяющее начальному условию  \eqref{vend:eq8} и нелокальным краевым условиям 
\[
D_\alpha\frac{\partial c}{\partial x}\Big\rvert_{x = 0}=D^{\alpha-1}_{0t}c(0;t),
\quad
\frac{\partial c}{\partial x}\Big\rvert_{x=l}=0.
\]
\end{newthm}

Для поставленной задачи доказана теорема существования и~единственности решения, предложен численно"=разностный алгоритм решения, реализованный модифицированным методом прогонки. 
Анализ численной реализации расчётного алгоритма, выполненный для различных значений параметра  $\alpha$, позволяет сделать вывод, что неустановившийся процесс загрязнения подземных вод с~течением времени стремится к~естественному 
равновесному состоянию, а~процесс распределения концентрации загрязнений имеет затухающий асимптотический характер, что является характерным для  нелинейного процесса переноса в~пористых средах.



\begin{thebibliography}{9}

\RBibitem{ven:bib:1}
\by Веригин~Н.\,Н.
\paper О кинетике растворения и выноса солей при фильтрации воды в грунтах
\inbook Растворение и выщелачивание горных пород
\publaddr М. 
\publ Госстройиздат
\yr 1957
\pages 17--27
\misctransl
\by Verigin~N.\,N.
\by On Kinetics of the Dissolution of Salts During the Filtration of Water into Soils
\inbook Dissolution and Leaching of Rocks 
\publ Gosstroyizdat
\publ Moscow 
\yr 1957
\pages 17--27



\RBibitem{ven:bib:2}
\by Вендина~А.\,А.
\paper Математическое моделирование массопереноса в~пористых средах
\jour Научная жизнь
\issue 3
\yr 2008
\pages 21--24
\misctransl 
\by Vendina~A.\,A.
\paper Mathematical modeling of mass transfer in porous mediums
\jour Nauchnaya zhizn'
\issue 3
\yr 2008
\pages 21--24




\RBibitem{ven:bib:3}
\by Нигматуллин~P.\,P.
\paper Дробный интеграл и~его физическая интерпретация
\jour ТМФ
\yr 1992
\vol 90
\issue 3
\pages 354--368
\transl англ. пер.:~
\paper Fractional integral and its physical interpretation
\by Nigmatullin~R.\,R.
\jour Theoret. and Math. Phys.
\yr 1992
\vol 90
\issue 3
\pages 242--251



\RBibitem{ven:bib:4}
\by Нахушев~А.\,М.
\book  Элементы дробного исчисления и их применение
\publaddr Нальчик
\publ КБНЦ РАН
\yr 2000
\totalpages 299
\misctransl
\by Nakhushev~A.\,M. 
\book Elements of Fractional Calculation and their Application
\publaddr Nal'chik
\publ KBNC RAN
\yr 2003
\totalpages 299

\end{thebibliography}


\noindent
\hskip6.5cm \thedate \par\noindent
\hskip6.5cm\therevdate

\vspace{-7mm}


%
%\chead[\fancyplain{}{\scriptsize \sl V\,e\,n\,d\,i\,n\,a~A.\,A.}]{\fancyplain{}{\scriptsize \sl  On the mathematical modeling of the contaminated groundwater fractal migration in natural porous systems}}

\makeengtitle

\newpage

\renewcommand{\baselinestretch}{0.9}

\newpage
%
%\end{document}